The reconstructed coordinates share inputs. In units with $c=1$, the implemented definitions (with reconstructed muon quantities) are
\begin{align}
 E_\nu^{\rm reco} &= E_\mu+q_0^{\rm reco},\nonumber\\
 Q^2_{\rm reco} &= \max\!\left\{0,\,
 2E_\nu^{\rm reco}(E_\mu-p_\mu\cos\theta_\mu)
 -m_\mu^2\right\},\nonumber\\
 q_3^{\rm reco} &= \sqrt{Q^2_{\rm reco}+(q_0^{\rm reco})^2},\nonumber\\
 W^{\rm reco} &= \sqrt{\max\{0,M^2+2Mq_0^{\rm reco}-Q^2_{\rm reco}\}},
 \label{eq:shared-reco-inputs}
\end{align}
where $m_\mu^2=E_\mu^2-p_\mu^2$ and $M$ is the average nucleon mass. The recoil-energy estimate $q_0^{\rm reco}$ differs from the tracker-plus-ECAL available energy $\Eavail^{\rm reco}$; truth $\Eavail$ also omits energy carried by species such as neutrons and is not the full energy transfer. Thus five recorded coordinates do not establish five independently constrained physical directions. Adding a coordinate, supplying a new detector input, and resolving a new truth direction are separate questions.
% Source: CVUniverse.h, NewEavail(), RecoQ3(), RecoW(), GetEAvailableTrue().
